% Adición a las conclusiones de VII; las conclusiones anteriores se conservan.
\paragraph{Integración de las cuatro interacciones y cierre canónico total.}
La gravitación y las interacciones fuerte, débil y electromagnética
actúan sobre una misma realización material de la estructura
APP--TRIT--TPK. El transporte conserva las hojas, la orientación,
la memoria y la incidencia del estado enriquecido; sus realizaciones
de espín y calibre se componen sobre las mismas fibras materiales.
La componente electromagnética pertenece a la realización
electrodébil: no se introduce como una interacción adicional
desconectada de ella. La acción total utiliza los coeficientes y
lectores ya construidos, incluidos la sección de Planck, la longitud
radial y el funcional de Barbero--Immirzi.

El contenido de esta integración no se limita a reunir términos en
una expresión. La variación de la misma acción produce la corriente
total de espín y la fuente métrica conjunta. La eliminación de la
contorsión conserva los términos cruzados entre especies; la
transformación de Legendre conserva los momentos materiales y
produce las restricciones totales. Las pruebas de las
secciones~\ref{hmtcuatro:gea:seccion}
y~\ref{hmtcuatro:accion-retroaccion-restricciones}
establecen la retroacción, la primera clase y la propagación de
esas restricciones en la realización regular especificada.

En particular, sea $\Theta_{\rm tot}$ el potencial canónico total
obtenido en esa transformación, $\mathcal Z=\{C_A=0\}$ la superficie
regular de restricciones y $X_N=N^AX_{C_A}$,
$X_M=M^BX_{C_B}$ dos deformaciones características admisibles.
La conexión inducida por la acción tiene
$H_N^{\rm tot,can}=-\Theta_{\rm tot}(X_N)$. Su curvatura, calculada
sin omitir las derivadas de los coeficientes de deformación, es
\begin{equation}
\begin{aligned}
\mathcal F^{\rm tot,can}_{N,M}
&=\delta_NH_M^{\rm tot,can}-\delta_MH_N^{\rm tot,can}
 -H_{[X_N,X_M]}^{\rm tot,can}
 +\frac{i}{\hbar}[H_N^{\rm tot,can},H_M^{\rm tot,can}]\\
&=-\mathrm d\Theta_{\rm tot}(X_N,X_M)
 =N^AM^B f_{AB}{}^{D}C_D.
\end{aligned}
\label{hmtcuatro:conclusion-curvatura-total}
\end{equation}
Por tanto,
\[
 \left.\mathcal F^{\rm tot,can}_{N,M}\right|_{\mathcal Z}=0.
\]
La fórmula incluye la deformación tangencial y las acciones de
calibre del corchete total; si se trabaja en el cociente, expresa
la misma igualdad módulo esas direcciones verticales. La
representación precuántica conserva este cierre mediante la
identidad de conmutadores demostrada en
\eqref{hmtcuatro:representacion-precuantica}.

Así, la incorporación gravitatoria es parte de la misma acción,
de sus fuentes y de su álgebra de restricciones, no una ecuación
yuxtapuesta después de construir las otras interacciones. La
anulación probada es la de la curvatura característica canónica;
no exige que sean nulas la curvatura del espacio-tiempo ni las
intensidades de los campos de calibre. Su dominio conserva el
coframe no degenerado, la carta regular y las condiciones de
borde utilizadas en la prueba. Tampoco suprime la memoria ni la
monodromía global de las historias que dieron lugar a la
realización.
